\documentclass[11pt]{article}

\usepackage[margin=1.05in]{geometry}
\usepackage{amsmath,amssymb,amsthm}
\usepackage{booktabs}
\usepackage{enumitem}
\usepackage{xcolor}
\usepackage[ruled,vlined,linesnumbered]{algorithm2e}
\usepackage[colorlinks=true,linkcolor=blue!45!black,citecolor=blue!45!black,urlcolor=blue!45!black]{hyperref}

\newcommand{\R}{\mathbb{R}}
\newcommand{\E}{\mathbb{E}}
\newcommand{\Var}{\operatorname{Var}}
\newcommand{\Cov}{\operatorname{Cov}}
\newcommand{\tr}{\operatorname{tr}}
\newcommand{\xmap}{x_{\star}}
\newcommand{\code}[1]{\texttt{\small #1}}
\newcommand{\Rhat}{\widehat{R}}

\theoremstyle{plain}
\newtheorem{proposition}{Proposition}
\theoremstyle{definition}
\newtheorem{condition}{Condition}
\theoremstyle{remark}
\newtheorem{remark}{Remark}

\title{\bfseries A fast, self-certifying Gaussian posterior\\[2pt] for MAGI ODE parameter inference}
\author{}
\date{\today}

\begin{document}
\maketitle

\begin{abstract}
\noindent
MAGI \cite{yang2021} infers ODE parameters from noisy, partially observed trajectories by placing a
Gaussian process prior on the state and conditioning on the ODE holding at a discretisation grid.
Its posterior is high dimensional --- the states outnumber the parameters by two orders of
magnitude --- so practitioners reach either for a long MCMC run or for the Laplace approximation
$\mathcal{N}(\xmap, H^{-1})$ at the maximum a posteriori point. We show the latter is worse than it
looks. On a FitzHugh--Nagumo benchmark it is mis-centred by a \emph{full posterior standard
deviation} on two of the three ODE parameters while its credible intervals have very nearly the
right width, so nothing in the output signals the error.

We give a pipeline that removes it in about one second. Three observations drive it. First, with
the observation noise fixed the MAGI log-density is \emph{exactly} a sum of squares, so the mode is
a nonlinear least-squares problem and Gauss--Newton reaches $\|\nabla \log p\| \approx 10^{-7}$ in
$0.05$\,s where a first-order optimiser stalls at $10^{3}$; the same structure yields the exact
Hessian from one Jacobian assembly rather than $d$ autodiff passes. Second, the residual error is
almost entirely in the mean, and the standard third-order Laplace correction removes $90$--$99\%$
of it. Third, a second and independently derived mean --- the stationary point of the reverse-KL
Gaussian objective, solved by a Newton iteration preconditioned on a screened low-dimensional
subspace --- supplies a \emph{reference-free} error estimate: the two derivations can only agree
closely if both are near the truth, and their disagreement relative to the size of the correction
separates the regimes where the expansion is valid from those where it is not by more than thirty
orders of magnitude.

We also give an exactness condition. If the vector field is affine in the state at fixed parameters
--- linear compartment models, linear pharmacokinetics, any constant-coefficient system, and much
weaker than joint affineness --- then $p(X \mid \theta)$ is exactly Gaussian, the $325$-dimensional
problem is exactly a $3$-dimensional one, and no approximation is required. The condition is
testable in two Hessian evaluations.

We report several negative results, including that Stein variational gradient descent does not
merely converge slowly here: started \emph{at} a correct answer it moves away from it, degrading
energy distance by a factor of $45$--$80$ while holding marginal standard deviations at $0.99$ of
the reference. Follow-up work reproduces that to three significant figures and locates the cause,
which is not specific to this posterior: the collapse is dimensional, and the bandwidth rule the
method ships with makes it a factor of $0.582K/\ln K$ worse than the rule it is usually credited
with. A corrected configuration exists and brings the ensemble to the reference's own agreement
floor---but it degrades the parameter estimates on two of three systems, which is what the exercise
is for, so we report it as a diagnosis rather than a recommendation.
\end{abstract}

\section{Introduction}

Inferring the parameters $\theta \in \R^{p}$ of an ODE system $\dot x = f(x,\theta)$ from noisy,
sparsely sampled observations is a classical inverse problem made awkward by the need to solve the
ODE inside the likelihood. MAGI \cite{yang2021} avoids the solver: it places a Gaussian process
prior on each state component, discretises time onto a grid $I = \{t_1,\dots,t_n\}$, and conditions
on the ODE residual being small at every grid point. The unknown becomes the concatenation of the
parameters and the state trajectory,
\[
  x = (\theta, X) \in \R^{d}, \qquad d = p + nD,
\]
with $D$ the number of state components. On the FitzHugh--Nagumo benchmark used throughout this
paper, $p=3$, $n=161$, $D=2$ and $d = 325$: the trajectory outnumbers the parameters $107$ to one.

This dimensionality is the practical difficulty. A well-tuned Hamiltonian sampler works but costs
minutes, and --- as we document in Section~\ref{sec:experiments} --- a configuration that converges
at one data density can fail to converge at another without that being obvious. The cheap
alternative is the Laplace approximation at the mode. Section~\ref{sec:diagnosis} shows that on
this problem the Laplace approximation is wrong in a specific and dangerous way: the covariance is
essentially right, so the credible intervals have the right \emph{width}, while the mean is off by
a full posterior standard deviation on two of the three parameters. An analyst reading the output
sees nothing amiss.

\paragraph{Contributions.}
\begin{enumerate}[leftmargin=1.5em,itemsep=1pt]
\item A characterisation of the Laplace error on MAGI posteriors: it lives in the mean, not the
covariance, and we localise the residual covariance error to a handful of directions that cancel
out of every marginal (Section~\ref{sec:diagnosis}).
\item A deterministic pipeline --- Gauss--Newton mode, exact structural Hessian, third-order mean
correction --- costing about one second at $d = 325$ and reducing the largest parameter error by
$2$ to $10\times$ at every setting tested (Sections~\ref{sec:lsq}--\ref{sec:correction}).
\item A reference-free certificate and the demonstration that a naive version of it fails: the
plain Gaussian-VI fixed point diverges, and restricting it to a screened subspace is what makes it
both convergent and cheap (Section~\ref{sec:certify}).
\item An exact discrepancy identity for the Gauss--Newton Gaussian, and an exactness condition on
the vector field under which no approximation is needed (Section~\ref{sec:theory}).
\item A methodological point that materially changed our own conclusions twice: in $d = 325$ the
usual covariance metrics are dominated by the reference chain's own sampling error and are
meaningless without a calibrated floor (Section~\ref{sec:floors}); and an error averaged over all
$d$ coordinates is dominated by trajectory states and can rank estimators backwards relative to the
parameters one actually wants (Section~\ref{sec:estimator}).
\end{enumerate}

\section{The MAGI posterior as nonlinear least squares}
\label{sec:lsq}

\subsection{Setup}

Write $\mu, \dot\mu \in \R^{n \times D}$ for the GP prior mean and its time derivative on the grid,
$C_d^{-1}, K_d^{-1} \in \R^{n\times n}$ for the prior and conditional-derivative precisions of
component $d$, $m_d$ for the corresponding conditional-mean operator, $\tau$ for the observation
mask, $y$ for the observations, $\sigma$ for the observation noise scale, and $\beta^{-1}$ for the
prior tempering factor. The MAGI log-density is
\begin{equation}
\label{eq:logp}
-2\log p(x) = \beta^{-1}\!\!\sum_d \|X_{\cdot d}-\mu_{\cdot d}\|^2_{C_d^{-1}}
  + \sum_{d}\frac{\|\tau_{\cdot d}\odot (X_{\cdot d}-y_{\cdot d})\|^2}{\sigma_d^2}
  + \beta^{-1}\!\!\sum_d \|r_d(x)\|^2_{K_d^{-1}} + c(\sigma),
\end{equation}
with ODE residual $r_d(x) = f_d(X,\theta) - \dot\mu_{\cdot d} - m_d (X_{\cdot d}-\mu_{\cdot d})$.

\subsection{Exact sum-of-squares structure}

\begin{proposition}[Least-squares form]
\label{prop:lsq}
With $\sigma$ held fixed, \eqref{eq:logp} is exactly a sum of squares. Writing $L_cL_c^\top =
C^{-1}$ and $L_kL_k^\top = K^{-1}$ and $b = \beta^{-1}$,
\begin{equation}
\label{eq:residual}
  -2\log p(x) = \|R(x)\|^2 + \text{const},\qquad
  R(x) = \begin{bmatrix}
    \sqrt{b}\, L_c^\top (X-\mu)\\[2pt]
    (X-y)/\sigma \;\text{ on observed entries}\\[2pt]
    \sqrt{b}\, L_k^\top\big(f(X,\theta)-\dot\mu-m(X-\mu)\big)
  \end{bmatrix} \in \R^{3nD}.
\end{equation}
\end{proposition}

\noindent
We verified \eqref{eq:residual} numerically against \eqref{eq:logp} to $8.4\times10^{-12}$. Two
consequences carry the rest of the paper.

\emph{Only the third block is nonlinear, and only through $f$.} Hence the residual Jacobian is
assembled from $n$ small \emph{pointwise} derivatives $\partial f/\partial(X_i,\theta) \in
\R^{D\times(D+p)}$ obtained by autodiff on the user's vector field, plus constant precomputed
blocks. The Gauss--Newton and observation blocks contribute a constant term to $J^\top J$, so the
normal equations are formed without ever materialising $J \in \R^{3nD \times d}$.

\emph{The second derivative of the residual is the second derivative of the ODE.} Differentiating
\eqref{eq:residual} twice,
\begin{equation}
\label{eq:hessR}
  \nabla^2 R_a = \sqrt{b}\,(L_k^\top)_a\, \nabla^2 f \qquad\text{(pointwise in the grid index)},
\end{equation}
so the exact Hessian of $U = -\log p$ is
\begin{equation}
\label{eq:hess}
  H(x) = \nabla^2 U = J(x)^\top J(x) + \textstyle\sum_a R_a(x)\,\nabla^2 R_a(x),
\end{equation}
in which the second term is $n$ small $(D{+}p)\times(D{+}p)$ blocks scattered into place. This
costs one Jacobian assembly rather than the $d$ forward passes \code{jax.hessian} would take; we
measured agreement with \code{jax.hessian} to $2.5\times10^{-9}$.

\subsection{The Gauss--Newton mode}

Proposition~\ref{prop:lsq} makes the mode a nonlinear least-squares problem, so Gauss--Newton
applies, damped in the manner of Levenberg and Marquardt, and converges quadratically near the
solution. This matters more than it may appear. The FitzHugh--Nagumo Hessian is conditioned around $2\times10^4$, where a
first-order optimiser stalls far from the mode \emph{while appearing to converge}: Prodigy
terminates at $\|\nabla\log p\| = 1.3\times 10^3$, a log-density $8.7$ nats below the mode, whereas
Gauss--Newton reaches $2.5\times10^{-7}$ in $0.05$\,s. Everything downstream --- the
metric, the correction, the certificate --- inherits that error. An earlier iteration of this work
was conducted on the stalled mode and reached materially different conclusions; the displacement
from mode to posterior mean was inflated threefold.

Unknown $\sigma$ is handled by block coordinate descent: Gauss--Newton on $(\theta, X)$ at fixed
$\sigma$, alternating with the exact conditional maximiser $\sigma_d^2 = \sum_{\text{obs}}(X-y)^2/N_d$.
Both blocks decrease the same objective. The whole iteration runs inside one \code{lax.while\_loop}
with no host synchronisation and one compilation per solver instance.

\section{Diagnosis: the error is in the mean}
\label{sec:diagnosis}

\subsection{Calibrating the metrics}
\label{sec:floors}

Before any comparison we must know what agreement is achievable. In $d = 325$ the affine-invariant
covariance metrics are dominated by the error in the \emph{reference's own} covariance estimate:
with $n_{\text{eff}}$ effective draws the eigenvalues of a sample covariance spread by roughly
$\pm 2\sqrt{d/n_{\text{eff}}}$, and a metric that is an rms over all $325$ of them inherits that.

We use three: the mean error $\text{bias} = \|\Sigma_g^{-1/2}(\mu_q-\mu_g)\|/\sqrt d$ in posterior
standard deviations per dimension; the F\"orstner distance
$\|\log(\Sigma_g^{-1/2}\Sigma_q\Sigma_g^{-1/2})\|_F/\sqrt d$; and $\mathrm{KL}(q\,\|\,\mathcal{N}(\mu_g,\Sigma_g))$.
Splitting an eight-chain, $64{,}000$-draw gold standard four-and-four gives F\"orstner $0.1538$;
inverting that against synthetic draws puts the chain's effective sample size at $27{,}791$
($43\%$), from which the floor for a \emph{noiseless} approximation scored against the full chain is
F\"orstner $0.0767$, bias $0.0042$, KL $0.48$. For scale, inflating the gold covariance uniformly by
$1.25\times$ scores F\"orstner $0.2231$ --- indistinguishable from what the Laplace covariance
scores. Comparisons reported without this calibration are not interpretable, and we found this
changed the sign of our reading of one experiment.

\subsection{What the Laplace approximation gets wrong}

Table~\ref{tab:diagnosis} scores $\mathcal{N}(\xmap, H^{-1})$ against the gold standard on the
quantities a practitioner reads.

\begin{table}[t]
\centering
\small
\begin{tabular}{lccc}
\toprule
quantity & $\mathcal{N}(\xmap,H^{-1})$ & $+$ third order & gold half (floor) \\
\midrule
$\theta_b$ mean error      & $\mathbf{-0.971}$ sd & $-0.010$ sd & $+0.017$ sd \\
$\theta_c$ mean error      & $\mathbf{+1.032}$ sd & $+0.093$ sd & $-0.008$ sd \\
$\theta_a$ mean error      & $+0.246$ sd & $+0.034$ sd & $+0.009$ sd \\
trajectory max mean error  & $1.015$ sd  & $0.115$ sd  & $0.042$ sd \\[2pt]
$\theta$ sd ratios         & $1.00/1.01/0.94$ & $0.99/1.01/0.94$ & $1.00/0.99/1.01$ \\
trajectory sd ratio (median) & $0.983$ & $0.985$ & $0.993$ \\
\bottomrule
\end{tabular}
\caption{The Laplace error at the baseline data density. The interval \emph{widths} are within a
few percent of a second independent reference run; two of the three parameter estimates are
centred a full standard deviation away.}
\label{tab:diagnosis}
\end{table}

This is the failure mode that matters in practice, because it is invisible. A user inspecting
posterior summaries sees plausible intervals of plausible width.

\subsection{Where the residual covariance error lives}

Three independent lines say the covariance is not worth chasing.

\emph{The residual error is real but hides in the null space of every marginal.} Scoring the
Laplace covariance against each half of the gold chain separately gives per-direction
log-variance errors correlated at $0.975$ (a pure-noise control gives $0.852$), so the error is
genuine. But it concentrates in about twenty directions holding $17\%$ of the posterior variance:
binning by variance share, the top five directions ($78.5\%$ of the variance) have rms log-ratio
$0.040$, the next twenty ($17.3\%$) have $0.245$, and the remaining three hundred ($4.2\%$) have
$0.03$. Those twenty are combinations that cancel out of every parameter marginal and every
pointwise trajectory band, which is why Table~\ref{tab:diagnosis} shows standard deviations within
$1$--$6\%$ while the aggregate F\"orstner distance sits at three times its floor.

\emph{At this data density, a method that demonstrably fixes marginal variances does not move the
joint covariance.} We
built Tierney--Kadane Laplace marginals \cite{tierney1986} along screened directions (Section~\ref{sec:screen})
using a bordered Gauss--Newton profile. For a unit direction $v$ the marginal of $z = v^\top(x-\mu)$
satisfies, to second order,
\[
  \log p_{\text{marg}}(z) = -U_{\text{prof}}(z) - \tfrac12\log\det H_\perp(z) + \text{const},
  \qquad U_{\text{prof}}(z) = \min\{U(x) : v^\top(x-\mu) = z\},
\]
and two identities make it cheap: the constrained Gauss--Newton step is the unconstrained step plus
a multiple of $A^{-1}v$, reusing the \emph{same} Cholesky factor, and for unit $v$,
$\det(N^\top H N) = \det(H)\,(v^\top H^{-1} v)$ for any orthonormal basis $N$ of $v^\perp$, so the
restricted determinant never requires the $(d{-}1)$-dimensional basis. At $0.31$\,s per direction it
cut the mean marginal-variance error from $4.68\%$ to $1.66\%$, fixing the two worst directions from
$+8.7\%$ to $+0.8\%$ and $+6.7\%$ to $+1.0\%$. Its effect on the joint covariance was F\"orstner
$0.2243 \to 0.2237$, for $13$\,s. We do not recommend it \emph{here}; Section~\ref{sec:covscaling}
revises this for harder settings, where the F\"orstner distance turns out to be the wrong thing to
have judged it on.

\emph{Sampled metrics cannot resolve the effect.} At $800$ draws the standard error of a
variance-weighted ratio is comparable to the $3$--$5\%$ effect being measured; it reported the
correction as a slight loss while the analytic comparison showed each corrected direction moving
decisively toward truth. Analytic approximations should be scored analytically.

\section{Correcting the mean}
\label{sec:correction}

\subsection{Third-order Laplace expansion}

The standard third-order correction to the posterior mean \cite{tierney1986,kass1990} is
\begin{equation}
\label{eq:third}
  \E[x] \;\approx\; \mu_3 \;:=\; \xmap - \tfrac12 H^{-1}\nabla \tr\!\big(H^{-1}\nabla^2 U\big),
\end{equation}
evaluated here with the exact Hessian \eqref{eq:hess}. It costs one gradient of a scalar built from
a Hessian --- $0.39$\,s at $d=325$ once compiled.

That \eqref{eq:third} should be needed at all is worth making concrete. Along the softest
eigendirection of $H$ the potential is not remotely quadratic: its relative departure from the
quadratic model is $1.85$ at one posterior standard deviation, $19.7$ at four and $114$ at sixteen,
while along stiff directions it is $0$ to four decimals. That direction carries roughly $20\%$ of
the total posterior variance.

\subsection{Gaussian variational inference, and why the obvious iteration diverges}

Minimising $\mathrm{KL}(q\|p)$ over Gaussians gives, by the Bonnet and Price theorems,
\begin{equation}
\label{eq:vi}
  \E_q[\nabla\log p] = 0, \qquad \Sigma^{-1} = \E_q[-\nabla^2\log p],
\end{equation}
so the Laplace approximation is precisely the zeroth iterate of a fixed point built from quantities
already available. The natural iteration $\mu \leftarrow \mu + \Sigma\,\E_q[\nabla\log p]$
\textbf{diverges}, for two reasons with one cause.

The mean step overshoots by about $1.6\times$ (measured step $0.2751$ against a correct $\approx
0.17$, with Monte Carlo error only $0.0025$) because the correct Newton preconditioner is
$\E_q[H]$, not $H(\xmap)$. These differ enormously: on the screened subspace we measure
$\E_q[H]/H(\xmap)$ between $6$ and $10$ at the baseline setting and $2121$ at the sparsest. The
average curvature of the potential over the posterior is an order of magnitude above its curvature
at the mode.

The covariance step explodes because a $325\times325$ matrix average over a few hundred samples has
$O(\sqrt{d/n})$ eigenvalue noise that inversion amplifies; one step took the trace ratio from
$1.007$ to $2.125$.

\begin{remark}[Antithetic pairing is load-bearing]
Our first estimator of $\E_q[\nabla\log p]$ produced noise four times the size of the signal. The
$O(\delta)$ term dominates the variance and is odd, so evaluating at $\mu\pm\delta$ cancels it
exactly. Spherical cubature is worse than useless here: its $2d$ nodes sit at $\sqrt d \approx 18$
posterior standard deviations along a single axis, deep in the region where the degree-six
polynomial dominates, and it diverged immediately.
\end{remark}

\subsection{A curvature screen}
\label{sec:screen}

Both failures are cured by refusing to estimate what the samples cannot support. Along eigenvector
$v_j$ of $H$, with unit equal to one posterior standard deviation $s_j = \lambda_j^{-1/2}$, define
\begin{equation}
\label{eq:screen}
  q_j \;=\; \frac14\sum_{t\in\{-2,-1,1,2\}}\left|\,\frac{U(\mu + t s_j v_j)}{t^2/2} - 1\,\right|,
\end{equation}
which is zero if and only if the slice is exactly quadratic. It costs \emph{four batched
log-density evaluations for all $d$ directions} ($0.15$\,s), ranks them at Spearman $0.760$ against
the true error, and its top five directions capture $71\%$ of the variance-weighted Laplace error
(top ten, $85\%$).

The screen is a good detector and a poor estimator: along the softest direction the slice is $185\%$
steeper than quadratic at one standard deviation while the true marginal is only $8\%$ too narrow,
the gap being the complement relaxing as one moves. We use it only to choose a subspace.

\subsection{Low-rank Gaussian VI}

Let $S$ index the $m$ largest $q_j$ and $V_S$ the corresponding eigenvectors. Restricted to $S$,
$\E_q[H]$ is an $m\times m$ block --- a well-determined estimate instead of an impossible one --- and
because $S$ is spanned by eigenvectors of $H$ the corrected curvature is block diagonal in $H$'s
eigenbasis and inverts in closed form:
\begin{equation}
\label{eq:woodbury}
  A = H + V_S\Delta V_S^\top,\qquad
  A^{-1} = \sum_{i\notin S}\frac{v_iv_i^\top}{\lambda_i} + V_S\big(\mathrm{diag}(\lambda_S)+\Delta\big)^{-1}V_S^\top .
\end{equation}
Only $V_S^\top H(x) V_S$ is ever needed, so $m$ Hessian-\emph{vector} products per sample replace a
$d\times d$ assembly. The iteration $\mu \leftarrow \mu + A^{-1}\E_q[\nabla\log p]$ then converges
in three steps, the step size falling $0.166 \to 0.0147 \to 0.0034 \to 0.0004$.

\subsection{Which mean to report}
\label{sec:estimator}

We obtain two corrected means: $\mu_3$ from \eqref{eq:third} and $\mu_{\text{VI}}$ from
\eqref{eq:woodbury}. Both are $O(\Lambda)$-accurate with different $O(\Lambda^2)$ residuals, so
their midpoint is a natural third candidate.

Scored as an error averaged over all $325$ coordinates the midpoint wins ($0.0234$ against $0.0268$
at baseline). Scored on the ODE parameters the ranking reverses, and by a wider margin: see
Table~\ref{tab:estimator}. The average is dominated by the $322$ trajectory states; the three
parameters are what the method exists to estimate. \textbf{We default to $\mu_3$}, giving up $13\%$
on the average to gain a factor of three on the parameters.

\begin{table}[t]
\centering
\small
\begin{tabular}{lcccc}
\toprule
& \multicolumn{3}{c}{largest $|\theta$ error$|$ (reference sd)} & all-coordinate \\
\cmidrule(lr){2-4}
setting & $\xmap$ & $\mu_3$ & midpoint & bias, $\mu_3$ vs mid \\
\midrule
baseline & $1.033$ & $\mathbf{0.099}$ & $0.297$ & $0.0268$ vs $\mathbf{0.0234}$\\
half     & $1.579$ & $\mathbf{0.164}$ & $0.616$ & $0.0934$ vs $\mathbf{0.0508}$\\
noisy    & $2.607$ & $1.337$ & $\mathbf{1.069}$ & $0.4370$ vs $\mathbf{0.1719}$\\
quarter  & $1.760$ & $\mathbf{0.446}$ & diverged & $0.3356$ vs --- \\
\bottomrule
\end{tabular}
\caption{The two metrics rank the estimators oppositely. We report $\mu_3$; the midpoint is
available for callers whose target is the trajectory.}
\label{tab:estimator}
\end{table}

\subsection{When the covariance does matter}
\label{sec:covscaling}

The preceding subsection is a statement about the baseline data density, and it does not
generalise. From the identity \eqref{eq:identity} the leading mean shift is driven by the cubic
term and is $O(\Lambda)$, whereas a covariance correction requires the quartic term or the square
of the cubic, both $O(\Lambda^{2})$. The covariance error is therefore negligible at small
$\Lambda$ and grows \emph{faster} than the mean error --- so the two must converge as a problem
becomes more non-Gaussian, and the question is only where the crossover lies.

Across our four settings, which span a factor of $200$ in measured non-Gaussianity, the ratio of
(mean error / its floor) to (covariance error / its floor) falls from $15.4$ at baseline to
$3.3$--$6.6$ at the harder settings, as predicted. In the units a user reads, the largest error in
a parameter standard deviation is

\begin{center}\small
\begin{tabular}{lcccc}
\toprule
 & baseline & half & noisy & quarter \\
$\kappa_S$ & $10.4$ & $45.8$ & $166.8$ & diverged \\
\midrule
largest $\theta$ sd error, $H^{-1}$ & $6.9\%$ & $17.7\%$ & $\mathbf{70.0\%}$ & $36.6\%$\\
after the profile correction & $\mathbf{0.9\%}$ & $\mathbf{6.4\%}$ & $\mathbf{22.0\%}$ & $\mathbf{11.6\%}$\\
reference half-vs-half floor & $1.8\%$ & $1.7\%$ & $3.6\%$ & $3.1\%$\\
\bottomrule
\end{tabular}
\end{center}

At $\sigma = 0.5$ the Laplace standard deviation for $\theta_b$ is $1.70$ times the reference's ---
a credible interval $70\%$ too wide, against a half-vs-half ratio of $1.0035$, so this is real. The
covariance is emphatically worth chasing there.

The tool that fixes it is the profile construction dismissed above, applied along the parameter
coordinate axes rather than along eigenvectors, with two repairs: the quadrature grid must be
bracketed on the \emph{profile's} own scale rather than the Laplace scale (at $\sigma = 0.5$ a
$\pm4.5$ Laplace-sd grid reaches $\pm7.6$ true standard deviations, where $H_\perp$ loses positive
definiteness and the Cholesky fails), and nodes at which $H_\perp$ is indefinite must be dropped
rather than poisoning the direction. It costs about $1.8$\,s per parameter and takes $\theta_b$
from $1.700$ to $0.960$ at the noisy setting and from $1.366$ to $1.033$ at the sparsest.

It is not uniformly safe: at the noisy setting it degraded $\theta_c$ from $1.8\%$ to $22\%$ error
while repairing $\theta_b$ from $70\%$ to $4\%$. It improves the worst case at every setting we
tried and should be run when $\kappa_S$ is large, but it is a repair for a Gaussian that is
already in trouble, not a routine step.

Our original dismissal of the method was an artefact of scoring it on the affine-invariant
F\"orstner distance, which weights all $325$ directions equally and is dominated by its own noise
floor. The parameter marginals are what a practitioner reads, and on those the method works.

\section{Certification without a reference}
\label{sec:certify}

\subsection{Certificates}

Everything above needs a reference chain to score. For deployment the question is whether the
method knows when it is failing. Four quantities, none requiring a reference:

\begin{itemize}[leftmargin=1.5em,itemsep=1pt]
\item $d_A$: $\|H_{XX}(\theta,X_1)-H_{XX}(\theta,X_2)\|/\|H_{XX}\|$ at two states sharing one
$\theta$; zero exactly under Condition~\ref{cond:A} below. Two Hessian evaluations.
\item $q_{\max}$: the largest screen value \eqref{eq:screen}. Four batched log-density evaluations.
\item $\kappa_S$: $\max_j \E_q[H]_{jj}/H_{jj}$ over the screened subspace.
\item $\rho$: the \emph{gate statistic} $\|\mu_3-\mu_{\text{VI}}\|_H/\|\mu_3-\xmap\|_H$.
\end{itemize}

$\rho$ is the important one. $\mu_3$ is a Taylor truncation of the posterior mean;
$\mu_{\text{VI}}$ is the exact stationary point of the reverse-KL Gaussian objective. Neither is the
truth and they come from genuinely different arguments, so they can only agree closely if both are
near it. Normalising by the size of the correction makes the statistic dimensionless and
self-scaling.

\subsection{What the gate does and does not diagnose}

Table~\ref{tab:certify} gives the certificates and the errors they are meant to predict.

\begin{table}[t]
\centering
\small
\begin{tabular}{lrrrr@{\hskip 1.2em}rrr}
\toprule
& \multicolumn{4}{c}{reference-free} & \multicolumn{3}{c}{needs a reference} \\
\cmidrule(lr){2-5}\cmidrule(lr){6-8}
setting & $q_{\max}$ & $\kappa_S$ & $\tau_{\text{end}}$ & $\rho$
        & $|\theta|$ err $\xmap$ & $|\theta|$ err $\mu_3$ & $\Rhat$ \\
\midrule
baseline (41 obs)   & $4.95$  & $10.4$    & $1{\cdot}10^{-4}$ & $0.198$ & $1.033$ & $0.099$ & $1.030$\\
half (21 obs)       & $20.8$  & $45.8$    & $3{\cdot}10^{-4}$ & $0.278$ & $1.579$ & $0.164$ & $1.121$\\
noisy ($\sigma$ 0.5)& $48.9$  & $167$     & $3{\cdot}10^{-2}$ & $0.244$ & $2.607$ & $1.337$ & $1.025$\\
quarter (11 obs)    & $932$   & $9.6{\cdot}10^{32}$ & $5{\cdot}10^{34}$ & $3.2{\cdot}10^{34}$ & $1.760$ & $0.446$ & $1.105$\\
\bottomrule
\end{tabular}
\caption{Certificates against truth. The gate statistic $\rho$ separates the regimes by more than
thirty orders of magnitude, but what it diagnoses is the breakdown of the VI solve, not the
validity of $\mu_3$.}
\label{tab:certify}
\end{table}

At the three denser settings $\rho$ lies in $0.20$--$0.28$; at the sparsest the VI iteration
diverges and $\rho$ exceeds $10^{34}$. Any threshold in that gap works; we use $0.5$.

We initially used $\rho$ to \emph{suppress} the correction, reporting $\xmap$ when the gate fired.
Better reference chains showed this to be wrong. At \code{quarter}, $\mu_3$ improves the largest
parameter error from $1.760$ to $0.446$ even though the VI has diverged; suppressing costs $1.3$
standard deviations of parameter accuracy to save about a fifth of a standard deviation averaged
over $322$ trajectory coordinates. The correction is therefore applied unconditionally and the gate
\emph{warns}: it reports that the low-rank VI has broken down and, with it, that the Gaussian over
the trajectory block should not be trusted, and that the user should escalate to a sampler.

\begin{remark}[A trust region destroys the statistic]
Bounding the VI step keeps a diverging iterate near $\mu_3$, which makes the two estimates look as
though they agree: at \code{quarter} a trust region pulled $\rho$ from $3\times10^{34}$ down to
$0.30$, below the threshold and wrong. The runaway magnitude is the signal and must not be damped.
\end{remark}

\section{Theory}
\label{sec:theory}

\subsection{An exact discrepancy identity}

Let $\delta = x - \xmap$ and let $E(\delta) = R(\xmap+\delta) - R_\star - J\delta$ be the residual's
nonlinear remainder, with $R_\star = R(\xmap)$.

\begin{proposition}[Exact discrepancy]
\label{prop:identity}
Let $q = \mathcal{N}(\xmap, (J^\top J)^{-1})$ be the Gauss--Newton Gaussian. Stationarity at the mode
gives $R_\star^\top J = 0$, whence, with no truncation of any kind,
\begin{equation}
\label{eq:identity}
  \log q(x) - \log p(x) \;=\; (R_\star + J\delta)^\top E(\delta) \;+\; \tfrac12\|E(\delta)\|^2 \;+\; \text{const}.
\end{equation}
\end{proposition}

\begin{proof}
$-2\log p = \|R_\star + J\delta + E\|^2 = \|R_\star\|^2 + 2R_\star^\top J\delta + \delta^\top J^\top
J\delta + 2(R_\star+J\delta)^\top E + \|E\|^2$. The cross term $R_\star^\top J\delta$ vanishes by
stationarity, and $-2\log q = \delta^\top J^\top J \delta + \text{const}$.
\end{proof}

By \eqref{eq:hessR} the entire departure from Gaussianity is thus carried by $\nabla^2 f$.

\subsection{An exactness condition}

\begin{condition}[Affine in the state]
\label{cond:A}
$f(\cdot,\theta)$ is affine in the state $X$ for every $\theta$.
\end{condition}

\begin{proposition}[Exact conditional Gaussianity]
\label{prop:condA}
Under Condition~\ref{cond:A} the residual is affine in $X$ at fixed $\theta$, $R = R_0(\theta) +
A(\theta)X$, so
\[
  U(\theta,X) = U\big(\theta,X^\star(\theta)\big) + \tfrac12(X-X^\star)^\top A^\top A (X-X^\star),
  \qquad X^\star(\theta) = -(A^\top A)^{-1}A^\top R_0,
\]
is exactly quadratic in $X$. Hence $p(X\mid\theta)$ is exactly Gaussian with covariance
$(A^\top A)^{-1}$, and integrating $X$ out is exact:
\begin{equation}
\label{eq:marginal}
  \log p(\theta) = -U\big(\theta, X^\star(\theta)\big) - \tfrac12\log\det\!\big(A(\theta)^\top A(\theta)\big) + \text{const}.
\end{equation}
\end{proposition}

The $d$-dimensional problem is then exactly a $p$-dimensional one, solvable by quadrature with no
MCMC and no Gaussian assumption on $\theta$; the only error is quadrature error, controlled by
adding nodes until the answer stops moving.

Condition~\ref{cond:A} is much weaker than requiring $f$ affine in $(X,\theta)$ jointly: $\theta$
may still multiply states, as it does in our test case through $c(V+R)$. It covers linear
compartment models, linear pharmacokinetics and every constant-coefficient system. Setting the
FitzHugh--Nagumo cubic coefficient to zero satisfies it, and we measure $d_A = 0.00\times10^{0}$
exactly there against $1.30\times10^{-1}$ for the true cubic field. Evaluating \eqref{eq:marginal}
with three structurally different quadrature rules --- Gauss--Hermite $9^3$, Gauss--Hermite $13^3$
and a uniform $21^3$ grid over $\pm5$ standard deviations --- agrees to all six printed decimals
(means $0.085448, 0.562092, 0.811322$), confirming convergence rather than a shared blind spot. The
Laplace parameter means there are off by $0.008$, $0.035$ and $0.027$ standard deviations.

\subsection{Why we do not report an a priori bound}

Identity \eqref{eq:identity} combined with $\|E(\delta)\| \le \tfrac12\Lambda\|\delta\|^2$,
$\Lambda = \sqrt{b}\,\|L_k\|\,L_2$ and $L_2 = \sup\|\nabla^2 f\|$, gives a rigorous bound. It is not
useful. Sweeping the cubic coefficient $\alpha$ from $0$ to $1$ we measure $\Lambda = 2012, 1588,
2548, 4173$ while the correction size varies as $0.005, 0.010, 0.111, 0.174$; the ratio
$\|\mu_3-\xmap\|/\Lambda$ varies by a factor of $17$. A supremum over the posterior bulk is far too
loose in $325$ dimensions. The measured certificates of Section~\ref{sec:certify} are the honest
substitute: they are computed rather than bounded, and they were validated against references.

The same sweep does show the diagnostics are coherent and monotone: at $\alpha = 0$, $d_A = 0$,
$\kappa_S = 1.18$ and $\rho = 0.0012$ --- every certificate independently and correctly reporting
that Laplace is exact --- rising monotonically to $d_A = 0.130$, $\kappa_S = 10.4$, $\rho = 0.198$ at
$\alpha = 1$.

\section{The pipeline}

Algorithm~\ref{alg:pipeline} assembles the pieces. Lines 1--3 are the deterministic answer and
account for two thirds of the cost; lines 4--13 exist to certify it, and produce a second mean as a
by-product. Nothing in the loop requires a $d\times d$ Hessian: line 10 needs only $m$
Hessian-vector products per sample, and line 11 inverts an $m\times m$ block through
\eqref{eq:woodbury}. The covariance is $H^{-1}$ throughout --- Section~\ref{sec:diagnosis} found
nothing cheap that improves on it.

Three choices in the algorithm are less arbitrary than they look, and each was arrived at by a
failure. The antithetic draw on line 8 is not variance reduction for its own sake but the
difference between a convergent and a divergent iteration. The spectral floor on line 10 is what
makes the Newton step a descent direction when sampling noise pushes an ill-conditioned block
indefinite, which happens in single precision. And the absence of a trust region on line 11 is
deliberate: bounding the step would make a diverging iterate resemble a converged one and destroy
the statistic computed on line 13.

\begin{algorithm}[t]
\DontPrintSemicolon
\SetAlgoLined
\KwIn{MAGI model; subspace size $m=12$; antithetic pairs $N=256$; gate $0.5$}
\KwOut{$\mathcal{N}(\mu,\Sigma)$ and certificates}
$\xmap \leftarrow$ Gauss--Newton on $\|R\|^2$ \tcp*[r]{$0.05$\,s}
$H \leftarrow J^\top J + \sum_a R_a\nabla^2R_a$ at $\xmap$; eigendecompose; $\Sigma \leftarrow H^{-1}$ \tcp*[r]{$0.30$\,s}
$\mu_3 \leftarrow \xmap - \tfrac12\Sigma\nabla\tr(\Sigma\,\nabla^2 U)$ \tcp*[r]{$0.39$\,s}
$q_j \leftarrow$ slice-curvature screen \eqref{eq:screen} for all $d$ directions \tcp*[r]{$0.15$\,s}
$S \leftarrow$ indices of the $m$ largest $q_j$\;
$\mu \leftarrow \xmap$\;
\For{$k = 1$ \KwTo $6$}{
  draw antithetic $\{\pm\delta_i\}_{i\le N}$ from $\mathcal{N}(0,\Sigma)$\;
  $g \leftarrow$ mean of $\nabla\log p$ over the $2N$ points\;
  $B \leftarrow$ mean of $V_S^\top H V_S$ by HVP; floor its spectrum\;
  $\mu \leftarrow \mu + A^{-1}g$ with $A^{-1}$ from \eqref{eq:woodbury} \tcp*[r]{no trust region}
}
$\rho \leftarrow \|\mu_3-\mu\|_H/\|\mu_3-\xmap\|_H$; warn if $\rho \ge 0.5$\;
\Return{$\mathcal{N}(\mu_3, \Sigma)$}
\caption{Deterministic MAGI posterior, $\approx 1.1$\,s at $d=325$}
\label{alg:pipeline}
\end{algorithm}

\section{Experiments}
\label{sec:experiments}

\subsection{Setup}

FitzHugh--Nagumo, $\dot V = c(V - V^3/3 + R)$, $\dot R = -(V - a + bR)/c$, on $t\in[0,20]$ with a
$161$-point grid, $d = 325$. Four settings vary the data: \emph{baseline} ($41$ observations per
component, $\sigma = 0.2$), \emph{half} ($21$), \emph{quarter} ($11$), \emph{noisy} ($41$, $\sigma =
0.5$). At \code{quarter}, $22$ data points constrain $325$ unknowns plus three parameters --- close
to unidentified, and not a regime in which any Gaussian approximation should be expected to work.

\subsection{References, and a warning about them}

References are NUTS \cite{hoffman2014} via BlackJAX, preconditioned by the exact Hessian: we sample
$y = H^{1/2}(x-\xmap)$ with an identity mass matrix. The metric affects only efficiency, never the
stationary distribution, so using the object under test as the preconditioner is not circular. The
baseline reference is $400$ chains of $150$ draws ($\Rhat = 1.03$); the three harder settings use
$64$ chains of $2000$, for reasons given below.

Credibility rests on the baseline, where an independent eight-chain gold standard exists: the
preconditioned run matches it \emph{better than the gold chain matches itself}, at bias $0.0056$
(half-vs-half floor $0.0095$), F\"orstner $0.1358$ (floor $0.1538$) and KL $1.50$ (floor $1.90$).

We stress a negative finding about the fallback. Our first references used $400$ chains of $150$
draws --- the same total cost --- and reached $\Rhat$ of $1.16$, $1.48$ and $1.77$ on the three
harder settings while giving $1.03$ at baseline. Being the exact method does not make a particular
run correct, and a configuration validated at one data density can silently fail at another. The
longer configuration reaches $\Rhat = 1.03$, $1.12$, $1.03$, $1.11$; \code{half} nonetheless reports
$6.2\%$ divergences, so we treat it as indicative. Conclusions in Table~\ref{tab:certify} that
depend on \code{half} and \code{quarter} should be read accordingly.

\subsection{Main result}

Table~\ref{tab:estimator} is the headline: the largest parameter error falls by factors of $10.4$,
$9.6$, $1.9$ and $3.9$ across the four settings. In energy distance over the whole vector, measured
in Mahalanobis coordinates against a floor set by scoring one $2000$-draw subsample of the reference
against another, the corrected Gaussian reaches $0.0397$ at baseline against a floor of $0.0331$ ---
at $2000$ draws it is not distinguishable from the reference chain.

\subsection{Ablations}

\emph{Hyperparameters.} $m = 12$ and $m = 24$ give \emph{identical} results, so the screened
subspace saturates --- an independent confirmation that the screen finds the right subspace --- while
$m=6$ is materially worse ($\rho = 0.064$ against $0.034$). $N$ from $256$ to $2048$ and three seeds
move $\rho$ by $\pm3\%$. The knobs were held fixed across every setting reported, so no result here
is the product of per-setting tuning.

\emph{Precision.} Single precision is the production default and is essentially free: the
correction agrees with its float64 value to $6.7\times10^{-4}$ relative, a gap of $0.00007$ posterior
standard deviations against a correction of $0.17$ and a floor of $0.006$, and the resulting bias is
identical to four decimals. The float32 mode stalls at $\|\nabla\log p\| = 2.9\times10^{-2}$ rather
than $5.4\times10^{-7}$, far below the level that affects anything downstream. This does depend on
forcing full-precision matmuls: the einsums in the log-density are matrix--vector contractions in
the forward pass, but their VJPs are matrix--matrix products, so on tensor-core hardware the
\emph{gradient} loses about four significant digits by default ($5.1\times10^{-7} \to 3.9\times10^{-3}$)
while the forward value is bit-identical. The value is not the thing to test. At the \code{noisy}
setting the VI solve does diverge in float32 where float64 converges; the gate catches it.

\subsection{Cost}

Algorithm~\ref{alg:pipeline} costs about $1.1$\,s at $d = 325$, against $144$\,s for the shortest
preconditioned NUTS run that is trustworthy at baseline and considerably more elsewhere. A cheaper
NUTS configuration ($400$ chains, $50$ warmup, $5$ draws, $159$ gradient evaluations, $21$\,s)
already reaches the sampling floor in energy distance, so the honest comparison is roughly
$20\times$, not $130\times$, if only distributional summaries are wanted --- but that configuration
provides no usable $\Rhat$.

\section{Negative results}

We record these because each cost us time and none is obvious in advance.

\emph{Stein variational gradient descent does not converge slowly here; it converges to the wrong
thing.} Starting mSVGD \emph{at} the corrected Gaussian --- a known-good answer --- and running it,
energy distance degrades from $0.085$ to $3.80$ (density-reweighted kernel \cite{huang2023}) or
$6.89$ (standard RBF), and the Stein-identity ratio falls from $1$ to $0.24$ and $0.019$. The
reweighted kernel is the more instructive failure: it holds marginal standard deviations at
$0.991$ of the reference --- visually perfect credible intervals --- while being far from the target.
Since the starting point was correct, the motion is unambiguously away from it, and no
initialisation strategy can help.

We have since reproduced that to three significant figures across a change of integrator and a
change of GP hyperparameter fit, and the explanation we gave for it was wrong. It is not a property
of the MAGI posterior: SVGD started at exact draws from a plain $\mathcal N(0, I_{325})$ collapses
identically. The equilibrium variance obeys $\mathrm{Var}_{\mathrm{SVGD}}/\mathrm{Var} = \ln K/d$,
because the median heuristic measures the bandwidth \emph{from the ensemble}, so contraction
tightens the bandwidth, which permits more contraction. Ba et al.\ \cite{ba2022} prove the
$\Theta(1/d)$ scaling for the plain median heuristic, with constant $(e-1)^{-1}K/d$; the extra
$1/\ln K$ that Liu and Wang recommend, and that our implementation used, costs a further
$0.582K/\ln K$ --- a factor of $39$ at $K = 400$. Deleting it is one line, and improves the energy
distance on these posteriors by $1.8$--$2.4\times$.

\emph{A configuration that works, and does not solve the problem we have.} Running SVGD in
coordinates whitened by $H^{-1}$ at the mode, at a \emph{fixed} bandwidth and a small fixed step,
has a genuine attractor---reached from starts a factor of four too narrow and a factor of four too
wide---and turns the shape error into a pure scale error. A scale error is one number, and the
Stein ratio measures it without a reference, so rescaling the ensemble by $1/\sqrt{R}$ completes
the correction: energy distance $0.57$ to $0.96$ times the $K$-particle floor on the three systems
with usable references, against $1.03$ to $1.65$ for the Gaussian draws it started from. On HIV it
beats $400$ exact draws on a Kolmogorov--Smirnov comparison.

It is still not a recommendation, and the reason is worth stating. The largest \emph{parameter}
error goes from $0.077$ to $0.242$ on FitzHugh--Nagumo and $0.071$ to $0.401$ on Lorenz, and the
rescaling does not repair it, because it corrects a scale and this is a location error. Only HIV
improves ($0.096 \to 0.067$). The ensemble metric improves because it is dominated by the $322$
trajectory states; the three numbers the exercise exists to estimate get worse. That is the same
trade this paper declines in Section~\ref{sec:estimator}, and it is declined here for the same
reason.

\emph{The Tierney--Kadane profile machinery is not worth it at the baseline density} --- correct,
$13$\,s, no gain on the joint covariance --- but this does not extend to harder settings, where it
is the only thing we found that repairs a badly wrong parameter interval
(Section~\ref{sec:covscaling}). We record it here mainly because judging it on the F\"orstner
distance was a mistake.

\emph{Unrestricted Gaussian VI diverges}, and a trust region that stops it diverging destroys the
diagnostic (Section~\ref{sec:certify}).

\emph{Spherical cubature} is unusable in this geometry: its nodes sit $18$ posterior standard
deviations out.

\emph{``Self-consistent Laplace''} --- recomputing $H$ at $\mu_3$ and reapplying \eqref{eq:third} ---
diverges by roughly nine orders of magnitude, because \eqref{eq:third} is derived \emph{at the
mode}, where $\nabla U = 0$; restoring the missing $-\Sigma\nabla U$ term simply recovers the VI
iteration.

\section{Limitations}

The evidence is one ODE system at four data densities. The gate threshold sits in a gap spanning
thirty orders of magnitude and is not delicate, but the claim that $\mu_3$ improves the parameters
\emph{unconditionally} rests on four settings and should be re-examined on other vector fields,
particularly stiff and multimodal ones. Two of the four reference chains are imperfect ($\Rhat =
1.12$ with $6.2\%$ divergences, and $\Rhat = 1.11$).

The pipeline leaves the covariance at $H^{-1}$, which Section~\ref{sec:covscaling} shows is
defensible only while $\kappa_S$ is small. We do not currently fold the profile repair into the
pipeline, because it is not uniformly safe and we have four settings on one vector field on which
to judge it; a user whose $\kappa_S$ is large should either run it per parameter and check, or use
a sampler.

Unknown observation noise is handled in the mode but the Gaussian approximation over $\sigma$ is
weaker than over $(\theta,X)$, since $\sigma$ is positivity-constrained and enters the log-density
through a clip; the structural Hessian \eqref{eq:hess} does not cover it and we fall back to a
generic one.

Finally, the pipeline returns a Gaussian. Where the posterior is genuinely non-Gaussian --- which the
certificates will say --- there is no substitute for a sampler.

\section{Conclusion}

The MAGI posterior is exactly a nonlinear least-squares problem, and taking that structure seriously
pays three times over: a mode accurate to $10^{-7}$ in $0.05$\,s, an exact Hessian from one Jacobian
assembly, and an explicit characterisation of what makes the posterior non-Gaussian. The
consequence of practical importance is that the default Laplace approximation is mis-centred by a
full posterior standard deviation on the parameters of interest while looking entirely healthy, and
that about a second of arithmetic fixes it and reports how far it can be trusted.

\begin{thebibliography}{9}
\small
\bibitem{yang2021} S.~Yang, S.~W.~K.~Wong, S.~C.~Kou.
  Inference of dynamic systems from noisy and sparse data via manifold-constrained Gaussian processes.
  \emph{PNAS} 118(15), 2021.
\bibitem{tierney1986} L.~Tierney, J.~B.~Kadane.
  Accurate approximations for posterior moments and marginal densities.
  \emph{JASA} 81(393):82--86, 1986.
\bibitem{kass1990} R.~E.~Kass, L.~Tierney, J.~B.~Kadane.
  The validity of posterior expansions based on Laplace's method. In
  \emph{Bayesian and Likelihood Methods in Statistics and Econometrics}, 1990.
\bibitem{hoffman2014} M.~D.~Hoffman, A.~Gelman.
  The No-U-Turn sampler. \emph{JMLR} 15:1593--1623, 2014.
\bibitem{ba2022} J.~Ba, M.~A.~Erdogdu, M.~Ghassemi, S.~Sun, T.~Suzuki, D.~Wu, T.~Zhang.
  Understanding the variance collapse of SVGD in high dimensions.
  \emph{ICLR}, 2022.
\bibitem{huang2023} C.~Huang, H.~Dong, W.~Fang.
  Detecting and mitigating mode-collapse for flow-based sampling / density-reweighted Stein kernels, 2023.
\bibitem{opper2009} M.~Opper, C.~Archambeau.
  The variational Gaussian approximation revisited. \emph{Neural Computation} 21(3):786--792, 2009.
\bibitem{forstner2003} W.~F\"orstner, B.~Moonen.
  A metric for covariance matrices. In \emph{Geodesy--The Challenge of the 3rd Millennium}, 2003.
\end{thebibliography}

\end{document}
